% Options for packages loaded elsewhere
\PassOptionsToPackage{unicode}{hyperref}
\PassOptionsToPackage{hyphens}{url}
\documentclass[
  ignorenonframetext,
]{beamer}
\newif\ifbibliography
\usepackage{pgfpages}
\setbeamertemplate{caption}[numbered]
\setbeamertemplate{caption label separator}{: }
\setbeamercolor{caption name}{fg=normal text.fg}
\beamertemplatenavigationsymbolsempty
% remove section numbering
\setbeamertemplate{part page}{
  \centering
  \begin{beamercolorbox}[sep=16pt,center]{part title}
    \usebeamerfont{part title}\insertpart\par
  \end{beamercolorbox}
}
\setbeamertemplate{section page}{
  \centering
  \begin{beamercolorbox}[sep=12pt,center]{section title}
    \usebeamerfont{section title}\insertsection\par
  \end{beamercolorbox}
}
\setbeamertemplate{subsection page}{
  \centering
  \begin{beamercolorbox}[sep=8pt,center]{subsection title}
    \usebeamerfont{subsection title}\insertsubsection\par
  \end{beamercolorbox}
}
% Prevent slide breaks in the middle of a paragraph
\widowpenalties 1 10000
\raggedbottom
\AtBeginPart{
  \frame{\partpage}
}
\AtBeginSection{
  \ifbibliography
  \else
    \frame{\sectionpage}
  \fi
}
\AtBeginSubsection{
  \frame{\subsectionpage}
}
\usepackage{iftex}
\ifPDFTeX
  \usepackage[T1]{fontenc}
  \usepackage[utf8]{inputenc}
  \usepackage{textcomp} % provide euro and other symbols
\else % if luatex or xetex
  \usepackage{unicode-math} % this also loads fontspec
  \defaultfontfeatures{Scale=MatchLowercase}
  \defaultfontfeatures[\rmfamily]{Ligatures=TeX,Scale=1}
\fi
\usepackage{lmodern}
\ifPDFTeX\else
  % xetex/luatex font selection
\fi
% Use upquote if available, for straight quotes in verbatim environments
\IfFileExists{upquote.sty}{\usepackage{upquote}}{}
\IfFileExists{microtype.sty}{% use microtype if available
  \usepackage[]{microtype}
  \UseMicrotypeSet[protrusion]{basicmath} % disable protrusion for tt fonts
}{}
\makeatletter
\@ifundefined{KOMAClassName}{% if non-KOMA class
  \IfFileExists{parskip.sty}{%
    \usepackage{parskip}
  }{% else
    \setlength{\parindent}{0pt}
    \setlength{\parskip}{6pt plus 2pt minus 1pt}}
}{% if KOMA class
  \KOMAoptions{parskip=half}}
\makeatother
\usepackage{color}
\usepackage{fancyvrb}
\newcommand{\VerbBar}{|}
\newcommand{\VERB}{\Verb[commandchars=\\\{\}]}
\DefineVerbatimEnvironment{Highlighting}{Verbatim}{commandchars=\\\{\}}
% Add ',fontsize=\small' for more characters per line
\newenvironment{Shaded}{}{}
\newcommand{\AlertTok}[1]{\textcolor[rgb]{1.00,0.00,0.00}{\textbf{#1}}}
\newcommand{\AnnotationTok}[1]{\textcolor[rgb]{0.38,0.63,0.69}{\textbf{\textit{#1}}}}
\newcommand{\AttributeTok}[1]{\textcolor[rgb]{0.49,0.56,0.16}{#1}}
\newcommand{\BaseNTok}[1]{\textcolor[rgb]{0.25,0.63,0.44}{#1}}
\newcommand{\BuiltInTok}[1]{\textcolor[rgb]{0.00,0.50,0.00}{#1}}
\newcommand{\CharTok}[1]{\textcolor[rgb]{0.25,0.44,0.63}{#1}}
\newcommand{\CommentTok}[1]{\textcolor[rgb]{0.38,0.63,0.69}{\textit{#1}}}
\newcommand{\CommentVarTok}[1]{\textcolor[rgb]{0.38,0.63,0.69}{\textbf{\textit{#1}}}}
\newcommand{\ConstantTok}[1]{\textcolor[rgb]{0.53,0.00,0.00}{#1}}
\newcommand{\ControlFlowTok}[1]{\textcolor[rgb]{0.00,0.44,0.13}{\textbf{#1}}}
\newcommand{\DataTypeTok}[1]{\textcolor[rgb]{0.56,0.13,0.00}{#1}}
\newcommand{\DecValTok}[1]{\textcolor[rgb]{0.25,0.63,0.44}{#1}}
\newcommand{\DocumentationTok}[1]{\textcolor[rgb]{0.73,0.13,0.13}{\textit{#1}}}
\newcommand{\ErrorTok}[1]{\textcolor[rgb]{1.00,0.00,0.00}{\textbf{#1}}}
\newcommand{\ExtensionTok}[1]{#1}
\newcommand{\FloatTok}[1]{\textcolor[rgb]{0.25,0.63,0.44}{#1}}
\newcommand{\FunctionTok}[1]{\textcolor[rgb]{0.02,0.16,0.49}{#1}}
\newcommand{\ImportTok}[1]{\textcolor[rgb]{0.00,0.50,0.00}{\textbf{#1}}}
\newcommand{\InformationTok}[1]{\textcolor[rgb]{0.38,0.63,0.69}{\textbf{\textit{#1}}}}
\newcommand{\KeywordTok}[1]{\textcolor[rgb]{0.00,0.44,0.13}{\textbf{#1}}}
\newcommand{\NormalTok}[1]{#1}
\newcommand{\OperatorTok}[1]{\textcolor[rgb]{0.40,0.40,0.40}{#1}}
\newcommand{\OtherTok}[1]{\textcolor[rgb]{0.00,0.44,0.13}{#1}}
\newcommand{\PreprocessorTok}[1]{\textcolor[rgb]{0.74,0.48,0.00}{#1}}
\newcommand{\RegionMarkerTok}[1]{#1}
\newcommand{\SpecialCharTok}[1]{\textcolor[rgb]{0.25,0.44,0.63}{#1}}
\newcommand{\SpecialStringTok}[1]{\textcolor[rgb]{0.73,0.40,0.53}{#1}}
\newcommand{\StringTok}[1]{\textcolor[rgb]{0.25,0.44,0.63}{#1}}
\newcommand{\VariableTok}[1]{\textcolor[rgb]{0.10,0.09,0.49}{#1}}
\newcommand{\VerbatimStringTok}[1]{\textcolor[rgb]{0.25,0.44,0.63}{#1}}
\newcommand{\WarningTok}[1]{\textcolor[rgb]{0.38,0.63,0.69}{\textbf{\textit{#1}}}}
\setlength{\emergencystretch}{3em} % prevent overfull lines
\providecommand{\tightlist}{%
  \setlength{\itemsep}{0pt}\setlength{\parskip}{0pt}}
\usepackage{bookmark}
\IfFileExists{xurl.sty}{\usepackage{xurl}}{} % add URL line breaks if available
\urlstyle{same}
\hypersetup{
  hidelinks,
  pdfcreator={LaTeX via pandoc}}

\author{\texorpdfstring{}{}}
\date{}

\begin{document}

\section{Confidence scoring, evidence tiers, and state-space
compilation}\label{confidence-scoring-evidence-tiers-and-state-space-compilation}

\begin{frame}[fragile]{Confidence Scoring and Evidence Tiers}
\protect\phantomsection\label{confidence-scoring-and-evidence-tiers}
The shipped \texttt{ConfidenceModel} in
\texttt{../cogant/py/cogant/translate/confidence.py} computes a scalar
score in \([0,1]\) from:

\begin{itemize}
\tightlist
\item
  Average confidence over provenance records.
\item
  An evidence-diversity term (scaled, capped).
\item
  A parser certainty factor applied multiplicatively.
\item
  Conflict penalties subtracted after scaling.
\end{itemize}

\begin{equation}
\label{eq:confidence-core}
c = \max\left(0,\ \min\left(1,\ (\bar{e} + \delta_d)\cdot \kappa - \pi\right)\right)
\end{equation}

Here \(\bar{e}\) is the mean evidence confidence, \(\delta_d\) is the
diversity bonus (bounded), \(\kappa\) is parser certainty, and \(\pi\)
aggregates conflict penalties. \textbf{Tiers} (for example static-only
versus static-plus-runtime) are assigned from score thresholds and
evidence source tags (\texttt{determine\_confidence\_tier}); see the
same module for named thresholds and enum values. The manuscript does
not duplicate those literals so they cannot drift from code.
\end{frame}

\begin{frame}[fragile]{State space and behavior}
\protect\phantomsection\label{state-space-and-behavior}
Where traces or coverage are available, \textbf{dynamic} extraction
feeds the state-space compiler. The goal is a compact behavioral model:
states, actions, transitions, and observations that sit alongside the
static graph for tasks that require execution-sensitive features. The
tuple \((S, A, T, O)\) of variables, actions, transitions, and
observation modalities mirrors the structure of a partially observed
Markov decision process as used in discrete active inference
{[}@parr2022active; @dacosta2020active; @smith2022stepbystep{]}, and at
the level of reachable state/transition graphs it also resembles the
Kripke structures traditionally used in model checking
{[}@clarke1999model{]}, without attempting to discharge temporal-logic
obligations. PyMDP {[}@heins2022pymdp{]} is one such downstream consumer
that executes the compiled state-spaces as discrete active inference
simulations over the exported Generalized Notation Notation bundles.

The shipped \texttt{StateSpaceCompiler} in
\texttt{../cogant/py/cogant/statespace/compiler.py} constructs this
model in several coordinated passes driven by the semantic mappings and
the underlying program graph.

\textbf{State variables} are identified by a
\texttt{StateVariableExtractor} that traverses graph nodes carrying
READS and WRITES edges: any node that participates in a write is a
candidate hidden-state variable, and its type, cardinality, and initial
confidence are derived from the node's recovered type string and the
provenance of the rules that produced the associated mapping.

\textbf{Actions} are extracted from semantic mappings of kind
\texttt{ACTION}, with parameters read from node metadata (typically the
parser-recovered function signature), effects traced through outgoing
WRITES edges on the controller node, and preconditions derived from
parameter lists, docstring directives, and explicit metadata entries.
For method-kind actions, the compiler also walks CONTAINS edges back to
the enclosing class so that instance-level state mutations are
attributed to the correct controller, mirroring how code property graphs
fuse structural and data-flow views {[}@yamaguchi2014modeling{]}.

\textbf{Transitions} are inferred by a cross-reference pass
(\texttt{\_cross\_reference\_actions\_and\_variables}) that, for each
action, partitions its adjacent variables into reads and writes based on
edge kind (WRITES and MUTATES yielding writes; READS and OBSERVES
yielding reads). The resulting \texttt{Transition} object records a
\texttt{source\_state} in which every touched variable is marked
\texttt{"pre"} and a \texttt{target\_state} in which written variables
advance to \texttt{"post"} while read-only variables remain
\texttt{"pre"}; this simple pre/post convention keeps the model faithful
to the static evidence without overcommitting to symbolic value domains
that cannot be recovered from AST analysis alone.

Trigger attribution follows incoming TRIGGERS and CALLS edges so that
orchestration flow is preserved in the behavioral model.

\textbf{Observation modalities} are built from semantic mappings of kind
\texttt{OBSERVATION}: each associated node becomes an
\texttt{ObservationModality} whose modality type (\texttt{log},
\texttt{metric}, \texttt{event}, \texttt{sensor}, or a generic fallback)
is inferred from the mapping's description and the node's name, giving
the GNN export a typed observation channel aligned with the OBSERVES
edges in the program graph.

The \textbf{temporal regime} of the model is determined by a companion
\texttt{TemporalAnalyzer}. It classifies nodes as asynchronous when
their metadata carries \texttt{is\_async}/\texttt{async} flags or their
names match patterns such as \texttt{async}, \texttt{callback},
\texttt{promise}, or \texttt{future}, and as event-related when their
kind is \texttt{EVENT} or their names match \texttt{event},
\texttt{handler}, \texttt{listener}, or \texttt{trigger}. Temporal
orderings are then extracted from CALLS and TRIGGERS edges, with each
edge classified as \texttt{parallel} (when either endpoint is async) or
\texttt{sequential} otherwise, and event patterns are assembled from
triggers-in / triggers-out pairs around each event node.

A final decision rule selects among \texttt{SYNCHRONOUS},
\texttt{ASYNCHRONOUS}, \texttt{EVENT\_DRIVEN}, and \texttt{HYBRID}: the
presence of event triggers and patterns combined with async handlers
yields \texttt{HYBRID}; event triggers alone yield
\texttt{EVENT\_DRIVEN}; an async fraction above 30 percent or the
presence of async handlers yields \texttt{ASYNCHRONOUS}; and all
remaining cases default to \texttt{SYNCHRONOUS}. This regime is attached
to the \texttt{StateSpaceModel} metadata so downstream consumers know
which execution model the transition graph assumes.

When coverage and traces are available, \texttt{enrich\_graph()} in
\texttt{../cogant/py/cogant/dynamic/enrichment.py} feeds additional
evidence into this pipeline. Coverage enrichment matches
\texttt{.coverage} SQLite databases or Cobertura XML reports against
nodes whose \texttt{path} and \texttt{source\_range} overlap covered
lines, attaching \texttt{coverage\_hits} and, where branch data is
available, \texttt{branch\_coverage} metadata. Trace enrichment parses
Chrome DevTools traces, writes \texttt{call\_count},
\texttt{avg\_duration\_ms}, and \texttt{is\_hot\_path} onto matching
callable nodes, and adds or reweights dynamic CALLS edges tagged with
\texttt{evidence\_sources={[}"dynamic\_trace"{]}}.

Both steps also append \texttt{dynamic\_coverage} and
\texttt{dynamic\_trace} markers to the program graph's evidence sources.
The confidence model consumes these markers directly: any mapping whose
evidence set now contains both static and dynamic entries becomes
eligible for promotion from the \texttt{STATIC\_ONLY} tier to
\texttt{STATIC\_PLUS\_RUNTIME}, and hot-path and branch-coverage
metadata raise the underlying score through the diversity bonus
\(\delta_d\) in Equation \ref{eq:confidence-core}. This is the mechanism
by which executing the target program, even partially, converts static
heuristics into corroborated behavioral facts without rerunning the
upstream rule engine {[}@allamanis2018survey{]}.

\begin{block}{Worked example: temperature controller}
\protect\phantomsection\label{worked-example-temperature-controller}
Consider a small Python controller extracted from a HVAC codebase:

\begin{Shaded}
\begin{Highlighting}[]
\KeywordTok{class}\NormalTok{ TemperatureController:}
    \KeywordTok{def} \FunctionTok{\_\_init\_\_}\NormalTok{(}\VariableTok{self}\NormalTok{):}
        \VariableTok{self}\NormalTok{.current\_temp: }\BuiltInTok{float} \OperatorTok{=} \FloatTok{20.0}
        \VariableTok{self}\NormalTok{.target\_temp: }\BuiltInTok{float} \OperatorTok{=} \FloatTok{22.0}
        \VariableTok{self}\NormalTok{.heater\_on: }\BuiltInTok{bool} \OperatorTok{=} \VariableTok{False}

    \KeywordTok{def}\NormalTok{ set\_target(}\VariableTok{self}\NormalTok{, t: }\BuiltInTok{float}\NormalTok{) }\OperatorTok{{-}\textgreater{}} \VariableTok{None}\NormalTok{:}
        \VariableTok{self}\NormalTok{.target\_temp }\OperatorTok{=}\NormalTok{ t}

    \KeywordTok{def}\NormalTok{ read\_sensor(}\VariableTok{self}\NormalTok{, reading: }\BuiltInTok{float}\NormalTok{) }\OperatorTok{{-}\textgreater{}} \VariableTok{None}\NormalTok{:}
        \VariableTok{self}\NormalTok{.current\_temp }\OperatorTok{=}\NormalTok{ reading}

    \KeywordTok{def}\NormalTok{ actuate\_heater(}\VariableTok{self}\NormalTok{) }\OperatorTok{{-}\textgreater{}} \VariableTok{None}\NormalTok{:}
        \ControlFlowTok{if} \VariableTok{self}\NormalTok{.current\_temp }\OperatorTok{\textless{}} \VariableTok{self}\NormalTok{.target\_temp:}
            \VariableTok{self}\NormalTok{.heater\_on }\OperatorTok{=} \VariableTok{True}
        \ControlFlowTok{else}\NormalTok{:}
            \VariableTok{self}\NormalTok{.heater\_on }\OperatorTok{=} \VariableTok{False}
\end{Highlighting}
\end{Shaded}

\texttt{StateVariableExtractor} identifies three \textbf{state
variables} from WRITES edges on \texttt{\_\_init\_\_} and the three
methods: \texttt{current\_temp} (float, cardinality continuous),
\texttt{target\_temp} (float), and \texttt{heater\_on} (bool,
cardinality 2). Three \textbf{actions} are extracted from
\texttt{ACTION}-kind mappings: \texttt{set\_target} (writes
\texttt{target\_temp}), \texttt{read\_sensor} (writes
\texttt{current\_temp}), and \texttt{actuate\_heater} (reads
\texttt{current\_temp}, \texttt{target\_temp}; writes
\texttt{heater\_on}). All three actions are attributed to
\texttt{TemperatureController} via CONTAINS edges.

The cross-reference pass yields three \textbf{transitions}. For
\texttt{actuate\_heater} the \texttt{source\_state} records
\texttt{\{current\_temp:\ "pre",\ target\_temp:\ "pre",\ heater\_on:\ "pre"\}}
and the \texttt{target\_state} records
\texttt{\{current\_temp:\ "pre",\ target\_temp:\ "pre",\ heater\_on:\ "post"\}},
capturing that only \texttt{heater\_on} advances while the read-only
variables remain pinned. Because no node carries async flags or
event-kind markers and no CALLS or TRIGGERS edges cross into async
endpoints, the \texttt{TemporalAnalyzer} classifies the model as
\texttt{SYNCHRONOUS} and attaches that regime to the
\texttt{StateSpaceModel} metadata.
\end{block}
\end{frame}

\end{document}
